\documentclass[10pt,a4paper]{amsart}
\usepackage[margin=19mm]{geometry}
\usepackage{amsmath,amssymb,mathtools}
\usepackage[scr=rsfs,cal=euler]{mathalfa}
\usepackage{xfrac}
\usepackage[unicode,hidelinks,bookmarks=false]{hyperref}
\newcommand{\ie}{i.e.\ }
\newcommand{\cf}{cf.\ }
\newcommand{\resp}{respectively\ }
\newcommand{\Subsection}{\subsection}
\providecommand{\nobreakdash}{\nobreak\hbox{-}}
\setlength{\emergencystretch}{2em}
\multlinegap0pt
\usepackage{bm}
\usepackage{bbm}
\usepackage{dsfont}
\usepackage{xspace}
\usepackage{cancel}
\usepackage{mathtools}
\usepackage{amsthm}
\makeatletter
\@ifundefined{thm}{\newtheorem{thm}{Thm}}{}
\@ifundefined{lem}{\newtheorem{lem}[thm]{Lemma}}{}
\makeatother
\makeatletter
\expandafter\ifx\csname com\endcsname\relax
\newenvironment{com}%
% begin comment
{\ifthenelse{\equal{\showcomments}{yes}}%
% then begin comment in margin
{\footnotemark
        \begin{lrbox}{\commentbox}
        \begin{minipage}[t]{1.5in}\raggedright\sffamily\tiny
        \footnotemark[\arabic{footnote}]}
% else eat contents of the environment
{\begin{lrbox}{\commentbox}}}%
% end comment
{\ifthenelse{\equal{\showcomments}{yes}}%
% then end comment
{\end{minipage}\end{lrbox}\marginpar{\usebox{\commentbox}}}
% else finish eating
{\end{lrbox}}}
\fi
\newcommand{\euler}{\chi}
\expandafter\ifx\csname refcom\endcsname\relax
\newenvironment{refcom}%
% begin comment
{\ifthenelse{\equal{\showrefcomments}{yes}}%
% then begin comment in margin
{\footnotemark
        \begin{lrbox}{\commentbox}
        \begin{minipage}[t]{1.25in}\raggedright\sffamily\tiny
        \footnotemark[\arabic{footnote}]}
% else eat contents of the environment
{\begin{lrbox}{\commentbox}}}%
% end comment
{\ifthenelse{\equal{\showrefcomments}{yes}}%
% then end comment
{\end{minipage}\end{lrbox}\marginpar{\usebox{\commentbox}}}
% else finish eating
{\end{lrbox}}}
\fi
\makeatother
\title[Cubical small-cancellation quotients of non-products]{Cubical small-cancellation quotients of non-products}
\author{Macarena Arenas, Kasia Jankiewicz, Daniel T. Wise}
\date{Source: arXiv:2509.13447v1 (CC BY 4.0)}
\begin{document}
\maketitle

\noindent\emph{Source and licence.} Cubical small-cancellation quotients of non-products. Version arXiv:2509.13447v1 (CC BY 4.0), distributed under the Creative Commons Attribution 4.0 International licence, as recorded in the arXiv licence metadata for this exact version and shown on the abstract page https://arxiv.org/abs/2509.13447v1. Full rights notice: licenses/arxiv-2509.13447-CC-BY-4.0.txt.\par\smallskip

\noindent\textit{Editorial scope.} Counted unit: the Lem whose source label is \texttt{no cover} in the article (source0.tex lines 657--659, source label \texttt{no cover}), delivered together with the article's own complete original proof of it (source0.tex lines 660--676, one \texttt{proof} environment of 17 lines). No mathematics is written, repaired or re-derived: statement and proof are the article's own text line for line. Required context retained uncounted: none. Every definition, display and label of those lines is retained unchanged. Omitted from this package and not redistributed: 1--656, 677--1806. Nothing inside a retained excerpt is omitted or altered: every retained body is delivered exactly as the article's lines are written, so no omission receipt of any kind accompanies this package. Citations: the retained excerpts carry no citation command at all, so no bibliography entry of the article is transcribed and no reference is redistributed. Reference adaptation: none was needed and none was made; every reference inside the delivered unit resolves inside it, so no unresolved reference or citation is produced. Printed slips kept literal: no printed word, symbol, hypothesis or number of the counted statement or of its proof was altered. Layout and class: the article's own class is \texttt{amsart} with packages including inputenc, babel, dutchcal, geometry, multicol, amsmath, amsthm, amssymb, graphicx, mathrsfs, todonotes, hyperref, comment, float; the wrapper is the lane's pinned \texttt{amsart} editorial wrapper at 10pt on A4, which restates the theorem-like environments the retained text needs and adds the article's own macro definitions that the retained text needs (\texttt{\textbackslash euler}), with extra wrapper packages (bm, bbm, dsfont, xspace, cancel, mathtools). The delivered numbering is the unit's own. Assets: no retained span contains a figure environment, a graphics-inclusion command, a PDF-page inclusion, a table or a TikZ picture; no image file is copied into this package, the package declares no asset dependency and no image file is redistributed with it. Licence: version arXiv:2509.13447v1 of this article is distributed under the Creative Commons Attribution 4.0 International licence, as recorded in the arXiv licence metadata for this exact version and shown on the abstract page https://arxiv.org/abs/2509.13447v1. A full rights notice is delivered at licenses/arxiv-2509.13447-CC-BY-4.0.txt. Characters: every retained excerpt is pure ASCII, so the delivered engine, XeLaTeX with its default Latin Modern text font, reports no missing character.
\begin{lem}\label{no cover}
    No  non-diagonal component of $Y\otimes_X Y$ is a finite cover of $Y$.
\end{lem}

\begin{proof}
Suppose $\widehat Y$ is a non-diagonal component of $Y\otimes_X Y$
that is a finite cover of $Y$.
Let  $x$ the basepoint of $X$, and let $y$ be a point of $Y$ mapping to $x$.
Let $(y,y')$ be a point of $\widehat Y$ that left-projects to $y$.
Since $(y,y')$ is not in the diagonal, we see that $y'\neq y$.
However, both $y$ and $y'$ project to $x$.
Let $\kappa$ be a path in $Y$ from $y$ to $y'$.
Then $[\kappa]\in \pi_1X$ stabilizes $\widetilde Y$,
but $[\kappa]\notin \pi_1Y$ since its based lift to $Y$ is not closed.
Thus $F\subsetneq \bar F= \langle F, [\kappa]\rangle$.
Let $\bar Y=\bar F \backslash \widetilde Y$.
Then there is a proper covering map $Y\rightarrow \bar Y$.
Let $d$ be the degree of this proper covering map, and note $1<d$ by properness
and $d<\infty$ since $Y$ is compact.
However, $-1 = \euler(Y)= d\cdot\euler(\bar Y) \neq -1$, which is impossible.
\end{proof}

\end{document}
